\section{Real Partitions And Squeezing}

The spline section shows how a finite set of anchors can force a balanced
completion. But a measurement grammar also has to live with refinement. A
partition can be split. A local support can be narrowed. A residual that was
invisible at one scale can become readable at another. If the chapter is to
reach a genuine invariant, it must explain how finite partitions approach a
completed reading without pretending that completion was available at the
start.

A real partition is a finite record laid across a real coordinate reading.
The word real does not license an infinite possession of every point. It
marks the kind of scale on which the finite record is being placed. The
grammar still carries only finitely many cells, finitely many anchors, and
finitely many supported tests. Each cell says: inside this bounded part of
the scale, these are the distinctions the ledger currently supports.

The finite truncation is therefore not a bad version of the continuum. It
is the only honest version the grammar can read at that stage. It carries
the supported data, the local completion rule, and the residual left by the
current mesh. A finer truncation may carry more data. It may split a cell,
add a knot, or test a narrower support. But it must preserve the older
commitments it has not explicitly revised. Refinement extends the ledger; it
does not give permission to forget why the older reading was admissible.

This is where the discrete Sobolev grammar enters. The grammar needs a way
to read not merely values, but activity: differences, bends, and supported
residuals. A Sobolev-style reading is the disciplined way of carrying both
the value information and the difference information. It measures whether a
finite record has bounded activity on its partition. It forbids a sequence
of refinements from hiding unbounded oscillation inside smaller and smaller
cells while pretending to converge as values.

The residual is then read weakly, against supported tests. A supported test
is a variation whose attention is confined to the part of the partition the
grammar has chosen. If every such test reads zero, the residual is invisible
on that supported ledger. If a test reads a nonzero value, the residual has
found a witness. The weak reading is weaker than omniscience and stronger
than silence: it asks exactly what the current partition is authorized to
ask.

Squeezing supplies the route to completion. At each finite stage, the
residual is bounded below by zero and above by a visible error bound. If the
upper bound is forced toward zero as the partition refines, the residual is
squeezed to zero. The completed reading is not asserted because someone has
looked at all points. It is read because every finite residual is trapped
between zero and a bound that vanishes under admissible refinement.

In symbols, the shape is simple:
\[
  0 \leq R_n \leq \epsilon_n,
  \qquad \epsilon_n \to 0 .
\]
The content is grammatical. The lower bound says no negative residue can be
used to cancel a hidden positive debt. The upper bound says the current
partition has measured the maximum unresolved imbalance it is still
carrying. The convergence of \(\epsilon_n\) says refinement is not wandering;
it is forcing the unread part to vanish. The limit is the reading permitted
by that forced vanishing.

The continuum-limit example is useful as a warning. Three coarse readings
can name a stress-strain relation, but they cannot determine every possible
continuum curve. A linear fit may converge to the wrong threshold if the
actual admissible shape contains a kink. Mesh refinement is honest only when
the residual bound being squeezed belongs to the grammar of the problem.
Refinement cannot rescue a model whose tests forbid the feature that the
record is trying to read.

This warning also protects the chapter from a common overreach. Completion
is not a license to add whatever smooth structure one likes. A completion is
admissible only when it is forced by finite partitions, supported tests, and
vanishing residual bounds. The Hilbert-style closure that later chapters
will use is therefore not a magical room added behind the finite record. It
is the completed home of the finite activity norm once the grammar has shown
that Cauchy refinements have a place to land.

Real partitions also explain why the spline result was not merely local
craftsmanship. A balanced cubic patch over one partition can be refined into
smaller patches. If the knot balances persist and the residual bound
shrinks, the finite completions form a tower. The tower carries values and
activity together. It reads convergence not by wishful smoothness but by the
disappearance of supported residuals.

Thus the grammar permits a continuum reading only at the end of a disciplined
squeeze. It starts with finite truncations, carries their supported
Sobolev activity, tests weak residuals, bounds the unread part, and forces
that bound to vanish. Completion is the quotient of all refinements that no
supported test can tell apart once the residual has vanished.

The next obstruction is subtler. A residual may vanish while a whole family
of invisible modes remains free. A constant drift, an affine tilt, or a
boundary-floating mode can carry no interior difference and yet still
prevent positivity. The grammar therefore needs an anchor.
